\section{Qualitative Analysis}
\label{sec:qual}

Table~\ref{tab:qualitative} shows three systems on two test videos that illustrate the omission/hallucination asymmetry; a no-crash example is in Table~\ref{tab:qualitative-full}.

\begin{table*}[!t]
\centering
\caption{Qualitative comparison on two held-out test videos. Explanations are lightly truncated; structured fields are the systems' own. Per-video ROUGE-L (R-L) against the reference Explanation is given for each output. \textcolor{red!70!black}{Red} marks a claim that contradicts the reference. A no-crash example (collision prior) is discussed in Section~\ref{sec:nocrash} and Table~\ref{tab:qualitative-full}.}
\label{tab:qualitative}
\renewcommand{\arraystretch}{1.12}
\scriptsize
\begin{tabular}{@{}p{1.5cm}p{0.42\textwidth}p{0.42\textwidth}@{}}
\toprule
 & \textbf{Video 001227} --- \emph{adaptation recovers the narrative} & \textbf{Video 000104} --- \emph{adaptation hallucinates the agent} \\
\midrule
\textbf{Reference} &
Curved two-lane road; a \emph{white car} approaches from the opposite direction; the camera car turns left to avoid it; front-left to front-left collision. Severity \code{moderate}; window $[3,4]$\,s.
&
Busy commercial road in daylight; a red bus passes; a slow \emph{black truck} ahead; the front of the camera car strikes the rear of the black truck. Severity \code{moderate}; window $[1,5]$\,s.
\\
\midrule
\textbf{Zero-shot Qwen2.5-VL-7B} &
``\#\#\# Analysis of the Car Crash Video --- Video Details: Date: 20/06/2014 --- Time: 17:25:23 --- Camera Perspective: \ldots'' No narrative reaches the crash. \textsc{No-Narrative}. R-L $0.00$.
&
``\ldots overcast weather conditions \ldots The road appears to be clear with no visible obstacles or hazards. Traffic seems'' (cut at 256 tokens). No collision described. \textsc{Fluent-Fabrication}. R-L $0.09$.
\\
\midrule
\textbf{\crashlogic{} (greedy)} &
Severity \textcolor{red!70!black}{severe} $|$ Start \textcolor{red!70!black}{2.0} End \textcolor{red!70!black}{3.0} $|$ Vehicles: White Car, Camera Car. ``\ldots a white car appeared from the opposite direction, speeding recklessly \ldots lost control \ldots impact between the front of the white car and the left side of the camera car.'' Agents and cause family match. R-L $0.45$, tIoU $0$.
&
Severity moderate $|$ Start 2.0 End 3.0 $|$ Vehicles: Camera Car, \textcolor{red!70!black}{Red Car}. ``\ldots the camera car followed closely behind a \textcolor{red!70!black}{red car} \ldots the red car abruptly applied its brakes \ldots the front of the camera car collided with the rear of the red car.'' Coherent rear-end story, wrong agent. \textsc{Wrong-Agent}. R-L $0.41$, tIoU $0.25$.
\\
\midrule
\textbf{\crashlogic{} + \tcd{}$_{0.5}$} &
Severity \textcolor{red!70!black}{severe} $|$ Start \textcolor{red!70!black}{2.0} End \textcolor{red!70!black}{3.0} $|$ Vehicles: Camera Car, White Car. ``\ldots a white car approached from the opposite direction on the left side \ldots both vehicles attempted to cross the intersection, resulting in a collision between the front-left of the camera car and the front-right of the white car.'' R-L $0.50$, tIoU $0$.
&
Severity moderate $|$ Start 2.0 End 3.0 $|$ Vehicles: Camera Car, \textcolor{red!70!black}{Red Car}. ``\ldots several cars were parked, while a few pedestrians were walking \ldots As the camera car attempted to overtake the black car in front of it, it encountered a \textcolor{red!70!black}{red car} that had come to a stop \ldots'' Same wrong agent; new unsupported details. R-L $0.31$, tIoU $0.25$.
\\
\bottomrule
\end{tabular}
\end{table*}


\paragraph{Video \code{001227}---adaptation recovers the narrative.}
The reference describes an oncoming white car on a curved two-lane road and a front-left to front-left impact after an evasive turn.
Zero-shot Qwen2.5-VL-7B emits markdown scaffolding and no narrative: \textsc{No-Narrative}, ROUGE-L $0$.
\crashlogic{} names the white oncoming car, a loss of control, and a front/left impact; agents and cause family match (ROUGE-L $0.45$).
\tcd{} keeps the agents and edges closer to the reference wording ($0.50$).
Both adapted outputs place the crash at $[2,3]$\,s; the reference window is $[3,4]$\,s (tIoU $0$).

\paragraph{Video \code{000104}---adaptation hallucinates the agent.}
The reference: the camera car strikes the rear of a slow \emph{black truck}; a red bus passes nearby.
Zero-shot text discusses weather and never reaches a collision.
\crashlogic{} constructs a coherent rear-end story---but with a \emph{red car} as the lead vehicle: \textsc{Wrong-Agent}.
\tcd{} keeps the red car and adds unsupported parked cars and pedestrians.
Under BLEU/ROUGE the adapted outputs score $0.41$ and $0.31$ against $0.09$ for zero-shot, so lexical metrics rank a specific wrong explanation far above a vague one.
This is the case that motivates scoring hallucination separately from omission.
